% ============================================================================
% BLOCK-003 Derivation v10: Tautology Audit + Order-of-Magnitude Check
% ============================================================================
\documentclass[11pt,a4paper]{article}

\usepackage{fontspec}
\usepackage{amsmath,amssymb,amsthm}
\usepackage{physics}
\usepackage{booktabs}
\usepackage{array}
\usepackage{longtable}
\usepackage{xcolor}
\usepackage[hidelinks]{hyperref}
\usepackage[margin=2.5cm]{geometry}
\usepackage{fancyhdr}
\usepackage{tcolorbox}

% Theorem environments
\newtheorem{proposition}{Proposition}
\newtheorem{lemma}{Lemma}

% Epistemic tags
\newcommand{\BL}{\textsc{[BL]}}
\newcommand{\Post}{\textsc{[P]}}
\newcommand{\Deriv}{\textsc{[Dc]}}
\newcommand{\Open}{\textsc{[Open]}}
\newcommand{\Taut}{\textcolor{red}{\textsc{[TAUT?]}}}

% Boxes
\newtcolorbox{auditbox}[1][]{colback=blue!5,colframe=blue!50!black,
  title=#1,fonttitle=\bfseries}
\newtcolorbox{resultbox}[1][]{colback=green!5,colframe=green!50!black,
  title=#1,fonttitle=\bfseries}

% Header/footer
\pagestyle{fancy}
\fancyhf{}
\fancyhead[L]{\small BLOCK-003 / Derivation v10}
\fancyhead[R]{\small EDC Gravity Program}
\fancyfoot[C]{\thepage}

\title{\textbf{EDC BLOCK-003 (Derivation v10):}\\[0.3em]
Tautology Audit and Order-of-Magnitude Sanity Check}
\author{EDC Working Document\\[0.5em]
\small 2026-02-02 --- Diagnostic before further attempts}
\date{}

\begin{document}
\maketitle

\begin{abstract}
After NC-1 (v8) and NC-2 (v9) both returned INCONCLUSIVE due to degenerate
parameter combinations, this note performs two diagnostics:
\textbf{Part A (Tautology Audit):} We trace where 4D gravity concepts
($G_N$, $1/r$ potential, Einstein equations) first enter the EDC framework
to assess whether any ``derivation'' of $G_N$ is circular.
\textbf{Part B (Order-of-Magnitude Check):} We set $\sigma = \sigma_0$
(a reasonable ansatz) and compute $G_N^{\rm pred}$ to check if EDC is
``in the right universe.''
\textbf{Results:} Tautology risk is MEDIUM---the 5D framework uses $\kappa_5$
without independent definition, potentially presupposing gravity. The
order-of-magnitude check shows EDC \textit{can} produce $G_N \sim 10^{-11}$
for reasonable $\sigma$, but requires fine-tuning.
\end{abstract}

\tableofcontents
\newpage

% ============================================================================
% PART A: TAUTOLOGY AUDIT
% ============================================================================
\part{Tautology Audit}

% ============================================================================
\section{Purpose}
% ============================================================================

A ``derivation'' of $G_N$ is \textbf{tautological} if 4D gravity was
already presupposed in the axioms or definitions. This audit traces the
logical chain to identify potential circularity.

\begin{auditbox}[Tautology Criterion]
A derivation is circular if:
\begin{enumerate}
\item $G_N$ (or Planck mass) appears as an input parameter.
\item The 4D Einstein equations are assumed before being ``derived.''
\item A $1/r$ Newtonian potential is used to define scales that later
``predict'' $1/r$ behavior.
\end{enumerate}
\end{auditbox}

% ============================================================================
\section{Assumption Ledger}
% ============================================================================

We trace where gravitational concepts enter the derivation chain.

\begin{longtable}{p{3cm}p{4cm}p{2cm}p{4cm}}
\toprule
\textbf{Concept} & \textbf{Where Introduced} & \textbf{Tag} & \textbf{Tautology Risk} \\
\midrule
\endhead
$\kappa_5^2$ (5D coupling) & 5D EH action $S_5 \sim \kappa_5^{-2} \int R^{(5)}$ &
\BL & \Taut\ if $\kappa_5$ defined via $G_5$ \\
$R^{(5)}$ (5D Ricci) & 5D metric variation & \BL & LOW (geometric) \\
$R^{(4)}$ (4D Ricci) & Induced on brane & \Deriv & LOW (follows from $R^{(5)}$) \\
$G_N$ & Matching $S_4 = \int R^{(4)}/(16\pi G_N)$ & \Deriv & MEDIUM (convention) \\
$1/r$ potential & 4D Newtonian limit & \Deriv & LOW if derived from $R^{(4)}$ \\
$M_{\rm Pl}$ & $M_{\rm Pl}^2 = 1/(8\pi G_N)$ & \BL & HIGH if used as input \\
$\sigma$ (brane tension) & Brane action $S_{\rm brane} = -\sigma \int \sqrt{-h}$ &
\BL & LOW (geometric) \\
$L$ (compact dim.) & Extra dimension size & \Post & LOW (geometric) \\
\bottomrule
\caption{Assumption ledger for gravitational concepts.}
\label{tab:ledger}
\end{longtable}

% ============================================================================
\section{Critical Analysis}
% ============================================================================

\subsection{The $\kappa_5$ Question}

The 5D gravitational coupling $\kappa_5$ is introduced in the bulk action:
\[
S_5 = \frac{1}{2\kappa_5^2} \int d^5X \sqrt{-g^{(5)}} R^{(5)}
\]

\textbf{Question:} How is $\kappa_5$ defined?

\begin{itemize}
\item \textbf{Option 1 (non-circular):} $\kappa_5$ is a fundamental constant
of the 5D theory, with dimensions $[\kappa_5^2] = L^3$, to be determined by
matching to observations.

\item \textbf{Option 2 (potentially circular):} $\kappa_5$ is defined via
$\kappa_5^2 = 8\pi G_5$, where $G_5$ is the ``5D Newton constant''---but
what defines $G_5$ if not by analogy to 4D gravity?
\end{itemize}

\textbf{In EDC:} $\kappa_5^2 = C \sigma^{-3/4}$ is proposed, where $C$ is
unfixed. This is \textit{not} directly circular (no $G_N$ input), but the
\textit{interpretation} of $\kappa_5$ as ``gravitational'' presupposes that
the 5D EH action describes gravity.

\subsection{The Matching Step}

The 4D Newton constant emerges from matching:
\[
\frac{L}{2\kappa_5^2} = \frac{1}{16\pi G_N}
\]

This is a \textit{definition} of $G_N$ in terms of $\kappa_5$ and $L$, not
a derivation. The question is whether $G_N$ so defined has the correct
physical interpretation (couples to stress-energy, produces $1/r$ potential).

\textbf{Assessment:} This step is LOW tautology risk---it is standard
dimensional reduction, not circular reasoning.

\subsection{The 4D Einstein Equations}

After reduction, the 4D Einstein equations emerge:
\[
G_{\mu\nu}^{(4)} = 8\pi G_N T_{\mu\nu}
\]

\textbf{Question:} Were these presupposed or derived?

\textbf{Answer:} Derived from the 5D equations + junction conditions. This
is legitimate \textit{if} the 5D equations are not themselves defined by
requiring 4D GR in some limit.

% ============================================================================
\section{Tautology Risk Assessment}
% ============================================================================

\begin{resultbox}[Tautology Risk: MEDIUM]
\textbf{Findings:}
\begin{itemize}
\item No explicit $G_N^{\rm obs}$ or $M_{\rm Pl}$ is used as input.
\item The 5D EH action is assumed (not derived from more fundamental principles).
\item The interpretation of $\kappa_5$ as ``gravitational'' is conventional,
not proven.
\item The matching step $G_N = \kappa_5^2/(8\pi L)$ is definitional, not
circular.
\end{itemize}

\textbf{Conclusion:} The derivation is not \textit{strictly} tautological,
but it assumes that 5D Einstein gravity is the correct bulk theory. This
is a \textit{model assumption}, not a circularity.

\textbf{Implication for BLOCK-003:} The problem is not tautology but
\textit{underdetermination}---the model has one free parameter ($C$ or
equivalent) that cannot be fixed without additional input.
\end{resultbox}

% ============================================================================
% PART B: ORDER-OF-MAGNITUDE CHECK
% ============================================================================
\newpage
\part{Order-of-Magnitude Sanity Check}

% ============================================================================
\section{Setup}
% ============================================================================

We set $\sigma = \sigma_0$ (an ansatz) and compute $G_N^{\rm pred}$ to check
if the order of magnitude is reasonable.

From v8:
\begin{equation}
G_N = \frac{C}{8\pi} \cdot \frac{\sigma^{-3/4}}{R_\xi}
\label{eq:gn_formula}
\end{equation}

We adopt the NP1 convention: $C = 8\pi$. Then:
\begin{equation}
G_N = \frac{\sigma^{-3/4}}{R_\xi} = \sigma^{-3/4} \cdot R_\xi^{-1}
\label{eq:gn_simple}
\end{equation}

% ============================================================================
\section{Natural Units and Conversion}
% ============================================================================

In natural units ($\hbar = c = 1$):
\begin{itemize}
\item $[G_N] = L^2 / M = L^2 \cdot L = L^3$ (since $[M] = L^{-1}$)
\item $[\sigma] = M / L^3 = L^{-4}$
\item $[R_\xi] = L$
\end{itemize}

Check Eq.~\eqref{eq:gn_simple}:
\[
[G_N] = [\sigma^{-3/4}][R_\xi^{-1}] = (L^{-4})^{-3/4} \cdot L^{-1}
= L^3 \cdot L^{-1} = L^2 \quad \text{???}
\]

\textbf{Problem:} Dimensions don't match! Let me recheck.

In SI units: $[G_N] = \text{m}^3 / (\text{kg} \cdot \text{s}^2)$.

In natural units where $\hbar = c = 1$ and we measure everything in GeV:
\begin{itemize}
\item $[G_N] = \text{GeV}^{-2}$
\item $[\sigma] = \text{GeV}^4$ (energy density)
\item $[R_\xi] = \text{GeV}^{-1}$
\end{itemize}

Then:
\[
[\sigma^{-3/4} / R_\xi] = \text{GeV}^{-3} / \text{GeV}^{-1} = \text{GeV}^{-2}
\quad \checkmark
\]

% ============================================================================
\section{Numerical Estimate}
% ============================================================================

\subsection{Input Values}

\begin{itemize}
\item $R_\xi = \hbar c / M_Z = 2.16 \times 10^{-18}$ m $\approx (91 \text{ GeV})^{-1}$
\item $G_N^{\rm obs} = 6.674 \times 10^{-11}$ m$^3$ kg$^{-1}$ s$^{-2}$
$\approx (1.22 \times 10^{19} \text{ GeV})^{-2} = M_{\rm Pl}^{-2}$
\end{itemize}

\subsection{Required $\sigma$ for Correct $G_N$}

From Eq.~\eqref{eq:gn_simple} with $C = 8\pi$:
\[
G_N = \sigma^{-3/4} R_\xi^{-1}
\]

Solving for $\sigma$:
\[
\sigma^{-3/4} = G_N \cdot R_\xi
\]
\[
\sigma = (G_N \cdot R_\xi)^{-4/3}
\]

In GeV units:
\begin{align}
G_N &\approx (1.22 \times 10^{19})^{-2} \text{ GeV}^{-2} \\
R_\xi &\approx (91)^{-1} \text{ GeV}^{-1}
\end{align}

\[
G_N \cdot R_\xi \approx \frac{1}{(1.22 \times 10^{19})^2 \times 91}
= \frac{1}{1.35 \times 10^{40}} \text{ GeV}^{-3}
\]

\[
\sigma = (1.35 \times 10^{40})^{4/3} \text{ GeV}^4
\approx 3.6 \times 10^{53} \text{ GeV}^4
\]

\subsection{Is This Reasonable?}

\begin{itemize}
\item $\sigma^{1/4} \approx 2.4 \times 10^{13}$ GeV $\approx 10^{-6} M_{\rm Pl}$
\item This is a GUT-scale energy density---high but not Planckian.
\end{itemize}

\textbf{Comparison:} In Randall-Sundrum, $\sigma \sim k M_5^3$ where $k$
is the AdS curvature. For $k \sim M_{\rm Pl}$, $\sigma \sim M_{\rm Pl}^4$,
which is much larger. The EDC value $\sigma \sim 10^{53}$ GeV$^4$ is
\textit{smaller} than Planckian, which is a mild hierarchy.

% ============================================================================
\section{Order-of-Magnitude Verdict}
% ============================================================================

\begin{resultbox}[Order-of-Magnitude: PLAUSIBLE]
\textbf{Finding:} To reproduce $G_N^{\rm obs}$ with $L = R_\xi$ and $C = 8\pi$,
the required brane tension is:
\[
\sigma \approx 3.6 \times 10^{53} \text{ GeV}^4
\quad \Leftrightarrow \quad
\sigma^{1/4} \approx 2 \times 10^{13} \text{ GeV}
\]

\textbf{Assessment:}
\begin{itemize}
\item This is GUT-scale, not Planck-scale---a mild hierarchy.
\item The model is ``in the right universe'' but requires $\sigma$ to be
specified to this precision.
\item No obvious contradiction with known physics.
\end{itemize}

\textbf{Caveat:} This is calibration, not prediction. We adjusted $\sigma$
to match $G_N^{\rm obs}$. True closure requires deriving $\sigma$ from
EDC field equations.
\end{resultbox}

% ============================================================================
\section{Combined Conclusions}
% ============================================================================

\begin{auditbox}[Summary]
\textbf{Part A (Tautology Audit):}
\begin{itemize}
\item Risk level: MEDIUM
\item No strict circularity found
\item Problem is underdetermination, not tautology
\end{itemize}

\textbf{Part B (Order-of-Magnitude):}
\begin{itemize}
\item EDC can produce correct $G_N$ for $\sigma \sim 10^{53}$ GeV$^4$
\item This is GUT-scale tension (plausible)
\item Requires $\sigma$ specification (same as before)
\end{itemize}

\textbf{BLOCK-003 Status:} Remains \Open.

The missing element is confirmed to be: \textbf{derive $\sigma$ from EDC
field equations} (or accept one calibration).
\end{auditbox}

\end{document}
